\documentclass[border=10pt]{standalone}
\usepackage{fontspec}
\AtBeginDocument{\hyphenation{zenodo}}% never hyphenate "zenodo" inside a DOI
\setmainfont{Noto Serif}[Ligatures=TeXOff]
\XeTeXgenerateactualtext=1
\newfontfamily\bengalifont{Noto Serif Bengali}[Script=Bengali, NFSSFamily=bnfam]
\usepackage[Latin,Bengali]{ucharclasses}
% group-free font switching (as in header-common.tex): robust inside braces, cells and quotes
\makeatletter
\def\BN@saved{\rmdefault}
\def\BN@bn{bnfam}
\newcommand\BNin{\edef\BN@tmp{\f@family}\ifx\BN@tmp\BN@bn\else\let\BN@saved\BN@tmp\fi\fontfamily{bnfam}\selectfont}
\newcommand\BNout{\fontfamily{\BN@saved}\selectfont}
\makeatother
\setTransitionsFor{Bengali}{\BNin}{\BNout}
% Typographic quotes and dashes are not in a ucharclasses class (class 0), so a Bangla word right
% after “ or ‘ would not trigger the Bangla font. Treat them as Latin so that every switch fires.
\XeTeXcharclass"2018=\XeTeXcharclass`A
\XeTeXcharclass"2019=\XeTeXcharclass`A
\XeTeXcharclass"201C=\XeTeXcharclass`A
\XeTeXcharclass"201D=\XeTeXcharclass`A
\XeTeXcharclass"2013=\XeTeXcharclass`A
\XeTeXcharclass"2014=\XeTeXcharclass`A
\XeTeXcharclass"2026=\XeTeXcharclass`A
\usepackage{newunicodechar}
\newunicodechar{।}{{\bengalifont\char"0964}}
\newunicodechar{॥}{{\bengalifont\char"0965}}

\usepackage[table]{xcolor}
\usepackage{booktabs,tabularx,array}
\definecolor{bapilrule}{HTML}{1F2937}
\definecolor{bapilstripe}{HTML}{F3F4F6}
\renewcommand{\arraystretch}{1.28}
\setlength{\tabcolsep}{5pt}

\begin{document}
\begin{minipage}{17cm}
{\large\bfseries সারণি ৩।}\ {\large হোয়াইট পেপারের অধ্যায় থেকে ব্লুপ্রিন্টের অধ্যায়}\par\vspace{7pt}
\small
\rowcolors{2}{bapilstripe}{white}
\begin{tabularx}{\linewidth}{@{}>{\hsize=0.892\hsize\raggedright\arraybackslash}X>{\hsize=0.649\hsize\raggedright\arraybackslash}X>{\hsize=1.459\hsize\raggedright\arraybackslash}X@{}}
\arrayrulecolor{bapilrule}\toprule
\rowcolor{white}\textbf{হোয়াইট পেপারের অধ্যায় (১.০)} & \textbf{ব্লুপ্রিন্টের অধ্যায়} & \textbf{ব্লুপ্রিন্ট যা যোগ করে} \\
\midrule
২. কেন বাংলাদেশের এআই নীতি-বুদ্ধিমত্তা প্রয়োজন & ২, ৪ & ভিত্তি-সূচক, ঘাটতি বিশ্লেষণ, ব্যবহারকারী-চরিত্র ও পরীক্ষাযোগ্য প্রয়োজনীয়তা \\
৩. বৈশ্বিক উদাহরণ & ৫, ১৩, ১৪ & দৃষ্টান্ত থেকে অভিযোজিত প্রাতিষ্ঠানিক পছন্দ ও নিয়ন্ত্রণ, কাঠামো-সংযোগ সারণিসহ \\
৪. বাংলাদেশে এআই-এর বর্তমান চিত্র & ২, ৩, ১৪ & পরিস্থিতি বিশ্লেষণ, অংশীজন নিবন্ধন ও ২০২৬ সালের আইনি কাঠামো \\
৫. BAPIL-এর রূপকল্প & ১, ১৬ & পরিসরের সীমারেখা, পরিপক্বতা শ্রেণি ও ২০২৬-২০৪১ কৌশলগত রোডম্যাপ \\
৬. BAPIL-এর স্থাপত্য & ৭, ৮, ১০ & বারোটি নামযুক্ত উপাদান, ইন্টারফেস, ডেটা-নিয়ম, হোস্টিং ও নিরাপত্তা অঞ্চল \\
৭. নীতি-বুদ্ধিমত্তা ড্যাশবোর্ড & ৭, ১২ & API বিনির্দেশ ও পাইলট বিনির্দেশ (P1, P2, P5) \\
৮. বাংলা নীতি ফাউন্ডেশন মডেল & ৯, ১২ & বাংলা টেক্সটের নিয়ম, BanglaPolicyQA বেঞ্চমার্ক ও প্রমাণনির্ভর মডেল-পথ \\
৯. এআই শাসন-কাঠামো & ১৩, সংযুক্তি খ, ঙ থেকে ট & ঝুঁকি-স্তর, জীবনচক্র-গেট, সনদ, কার্যপরিধি ও চেকলিস্ট \\
১০. জাতীয় নিরাপত্তা ও সাইবার নিরাপত্তা & ১০ & NIST CSF 2.0 সংযোগ, এআই-নির্দিষ্ট হুমকি, ঘটনার স্তর ও পুনরুদ্ধার-লক্ষ্য \\
১১. এআই শিক্ষা ও মেধা & ৬ & জনবল পরিকল্পনা, দক্ষতা-কাঠামো ও প্রশিক্ষণ পাঠ্যক্রম \\
১২. স্টার্টআপ ইকোসিস্টেম & ৩, ১৭ & শিল্পখাতের সম্পৃক্ততা ও বাজার উন্মুক্ত রাখার ক্রয়-শর্ত \\
১৩. সম্পদ ও অর্থায়ন & ১৭ & ব্যয়-মডেল, দিগন্তভিত্তিক অর্থায়ন-মিশ্রণ ও টেকসইতা পরিকল্পনা \\
১৪. পাঁচ বছরের রোডম্যাপ & ১৫ & মাইলফলক, ডেলিভারেবল ও গেটসহ ৩৬ মাসের কর্মপরিকল্পনা \\
১৫. দীর্ঘমেয়াদি রূপকল্প & ১৬ & ২০৪১ ও তার পরবর্তী সময় পর্যন্ত দিগন্ত, নির্দেশক লক্ষ্য, পরিস্থিতি ও পর্যালোচনা-সূচক \\
১৬. ঝুঁকি, নৈতিকতা ও দায়িত্বশীল এআই & ১৩, ১৯ & ঝুঁকি নিবন্ধন, তাপ-মানচিত্র ও পরিচালনাগত নিশ্চয়তা \\
১৭. কৌশলগত সুপারিশ & ২১ & কর্তৃপক্ষের জন্য ছয়টি সুনির্দিষ্ট সিদ্ধান্ত \\
পরিশিষ্ট ক. পাইলট প্রকল্প & ১২ & ব্যবহারকারী, ডেটা, স্তর, পরিসর ও উত্তরণ-মানদণ্ডসহ পূর্ণ বিনির্দেশ \\
\bottomrule
\end{tabularx}
\par\vspace{3pt}{\scriptsize\color{gray} উৎস: BAPIL টেকনিক্যাল ব্লুপ্রিন্ট, সংস্করণ ১.০ (২০২৬)। CC BY 4.0।}
\end{minipage}
\end{document}
